\section{Methods, Evidence, and Comparator Discipline}
\label{sec:oc133-methods-evidence-comparators}

This chapter explains how Core 1.3.3 asks to be reviewed. It is not a dump of
the evidence register. The public claim is always read through four questions:
what object is being named, what support class is actually present, what would
reopen the wording, and which comparator could explain the same observation
with less theoretical burden. The value of this discipline is that a reviewer
can attack a precise relation instead of arguing against a slogan.

\subsection{Claim wording and support class}

A sentence in the manuscript may be definitional, theorem-bound, finite-witness
bound, replay-bound, comparator-bound, illustrative, or planned. These classes
are intentionally different. A definition may introduce vocabulary without
proving empirical reach. A theorem-bound sentence may be strong inside stated
assumptions but silent outside them. A replay-bound sentence can support a
case row without becoming a completed domain validation. When the support class
changes, the public wording must change with it.

\subsection{The canonical model object}

The model object inspected across this release is
\[
  K=(\Omega,\partial\Omega,A,\Theta,P,J,C,k,M).
\]
The tuple is not merely notation. It gives a reviewer nine inspection points:
the admissible state region, the boundary, observable axes, threshold
conditions, potentials, flows, cycles, continuumness, and embedding context.
The article, monograph, figures, and tables may use prose glosses for these
components, but they do not promote an alternate canonical tuple.

\subsection{Formal and finite-witness support}

Formal support is cited only where the present release has a theorem row, proof
sheet, finite semantic witness, or formalization inventory entry that matches
the public statement. The current public boundary is deliberately cautious:
formalization inventory can orient the reviewer, but it is not described as a
clean machine-checked certificate unless the binding is clean for the exact
claim. If a proof sheet supports only a finite witness or a stated assumption
class, the public sentence remains inside that class.

\subsection{Replay and numeric support}

Replay rows are treated as bounded scientific checks. A row becomes
reader-facing support only when it names a source snapshot, reconstruction rule
or formula, observed or derived value, residual or uncertainty, comparator,
negative control, and falsifier. Missing pieces do not disappear; they demote
the row to an illustrative example or a planned measurement. The principal
prediction-row anchor is
\path{validation/target_blind/OC133_TARGET_BLIND_PREDICTION_TABLE.json}.

\subsection{Comparator discipline}

Prior systems theory, cybernetics, autopoiesis, dynamical systems, complexity
science, formal verification, formal ontology, category-theoretic modeling, and
reproducibility research are not treated as background scenery. They are active
comparators. The relevant public comparator anchors are
\path{docs/OC_1_3_3_PRIOR_ART_COMPARATOR_MATRIX.json} and
\path{comparators/OC_1_3_3_COMPARATOR_MATRIX.md}. If a comparator explains a
row with less burden, the OC claim narrows to its residual contribution.

\subsection{Reopening conditions}

Every promoted claim should be reopenable. A tuple claim reopens if a component
can be removed without changing the verdict. A K-level claim reopens if the
level adds no retained witness or demotion criterion. A proof claim reopens if
the dependency or assumption class is mismatched. A replay claim reopens if the
source snapshot, formula, comparator, residual, negative control, or falsifier
is missing or fails. This is not a weakness of the release; it is the mechanism
that keeps strong prose from outrunning visible evidence.

\subsection{Lifecycle and identity claims}

Lifecycle claims are especially easy to overread, so this release separates
death, residue, rebirth, and identity continuation. A residue is not a live
continuation merely because it carries traces of the source. A rebirth relation
is not same-identity survival unless endpoint-bound identity evidence is
explicitly present. The reviewer should therefore ask which token continues,
which invariant is declared, which endpoint evidence is available, and which
condition would make the identity claim false.

\subsection{K-level claims}

K-levels are witness disciplines rather than labels. A level must add an
observable or formal obligation: a retained witness, a constraint relation, a
demotion criterion, or a finite row that would fail if the level were merely a
name. The K0--K12 hierarchy is therefore inspected as a nested system grammar:
lower continua compose the next level, while higher contexts constrain the
admissible behaviour of lower ones.

\subsection{Article projection boundary}

The compact journal article is a projection from the release, not a parallel
theory. It may shorten examples and omit supporting appendices, but it may not
introduce unsupported scientific claims. Its four claims are the tuple, the
K-level witness discipline, the evidence-promotion rule, and the bounded
replay/comparator route. Each claim has a reopening condition printed in the
article so that the article can be criticized without treating the monograph as
a black box.

\subsection{Public reading rule}

The public reading rule is simple: do not infer broader closure from an anchor
than the anchor can carry. The release is strongest when it states exactly what
is definitional, what is theorem-bound, what is finite-witness-bound, what is a
bounded replay, and what remains future work. This rule is the bridge between
ambition and reviewability.

\subsection{Worked review example: tuple redundancy}

Suppose a critic argues that the boundary component can be absorbed into the
state region. The manuscript should not answer with rhetoric. It should ask
whether the case still distinguishes an admissible live state from a state that
forces demotion, death, split, or narrower wording. If that distinction remains
visible after removing \(\partial\Omega\), then the tuple row is too broad. If
the distinction disappears, the boundary component has earned local work in
that case. The same test applies to thresholds, flows, cycles, and embedding
context. Redundancy is not settled by how elegant the tuple looks; it is settled
by whether a component changes the review verdict.

\subsection{Worked review example: K-level promotion}

Suppose a section names a system as K7 because people, institutions, or
authority are involved. That is not enough. K7 promotion requires a retained
witness of social coordination: authority, trust, obligation, delegation,
norm-enforcement, or another declared coordination relation that cannot be
reduced to a lower-level component without loss in the stated case. If the row
only says that humans are present, it is an example candidate, not a promoted
K-level claim. If it also states the coordination witness, the lower-level
composition, the higher-level constraint, and the demotion condition, the row
can be reviewed as a K7 statement.

\subsection{Worked review example: replay support}

Suppose a retrospective bounded replay row reports a bounded match. The public conclusion
is still not ``the domain is validated.'' The conclusion is that the row, under
its source snapshot and reconstruction rule, did not trigger the named
falsifier and did not lose immediately to the declared comparator. The reviewer
then asks whether the case is representative, whether the residual is small for
the right reason, whether a negative control would fail, and whether a rival
model explains the row with less burden. Only after those questions survive can
the wording move beyond illustrative support.

\subsection{Worked review example: comparator pressure}

Suppose a cybernetic or dynamical-systems account already explains feedback,
state transitions, and attractor behaviour in a case. OC does not earn novelty
by renaming those concepts. It earns residual value only if the continuum tuple,
K-level witness discipline, lifecycle boundary, or evidence-promotion rule
adds an inspection point that the comparator does not already supply for the
same row. If no residual point remains, the honest action is to credit the
comparator and narrow the OC sentence.

\subsection{Worked review example: figure and table claims}

A diagram may make a claim by implication even when the caption is cautious. A
K-level figure, for example, can imply a complete hierarchy; a replay table can
imply empirical validation; a comparator table can imply novelty. In this
release a visual or table is treated as public argument only if its caption
states what it demonstrates, which variables or numbers matter, which formula
or evidence row supports it, and what would reopen it. A beautiful picture that
cannot answer those questions remains decoration and should not carry a claim.

\subsection{How to read absence}

Absence is also evidence. If a source snapshot is absent, the claim is not
inspectable. If a proof dependency is absent, the theorem wording is not
closed. If a comparator row is absent, novelty is not yet bounded. If a
falsifier is absent, the reader cannot tell what would make the claim false.
The manuscript therefore treats missing anchors as demotion signals rather than
as blank spaces to be filled by confidence.

\subsection{Claim-classification checklist}

The following checklist is the practical reading instrument for this chapter.
It is written for a reviewer who wants to decide quickly whether a paragraph is
claiming more than the release can support.

\begin{enumerate}
\item Is the sentence a definition, a theorem-bound statement, a finite-witness
      statement, a replay statement, a comparator statement, an illustration,
      or a planned research statement?
\item Does the sentence name or imply the canonical object
      \(K=(\Omega,\partial\Omega,A,\Theta,P,J,C,k,M)\), a K-level, a lifecycle
      relation, a proof relation, a replay row, a comparator row, or a journal
      projection?
\item Is the cited support class sufficient for the verb being used? Words such
      as proves, validates, explains, supports, illustrates, bounds, and
      suggests do not carry the same scientific force.
\item Does the paragraph state the reopening condition, or is the reader forced
      to infer it?
\item If the paragraph were moved into a journal article without the monograph,
      would the claim still be inspectable?
\end{enumerate}

If the answer to the final question is no, the correct repair is not merely a
footnote. The public sentence should be narrowed, the source anchor should be
made visible, or the material should be moved to an appendix where its support
class can be explained without pretending to be a compact journal result.

\subsection{Evidence verbs}

The manuscript uses evidence verbs conservatively. A definition \emph{names} a
distinction. A theorem-bound row \emph{establishes} a result only under its
assumptions. A finite witness \emph{shows consistency or counterexample
behaviour} in the finite case. A replay row \emph{supports} a bounded reading
only when its source snapshot and reconstruction rule are visible. A comparator
row \emph{bounds} novelty; it does not make novelty automatic. An illustrative
example \emph{orients} a reader; it does not validate a domain. This verb
discipline is intentionally plain because it is one of the easiest ways to
prevent overclaim.

\subsection{Promotion and demotion}

Promotion is reversible. A claim may move from planned to illustrative, from
illustrative to replay-bound, from replay-bound to stronger domain evidence, or
from theorem-oriented prose to formal support as the corpus improves. It can
also move downward. If a source row is missing, a replay row fails, a comparator
absorbs the residual contribution, or a proof dependency does not match the
public wording, demotion is the honest scientific action. The release is
therefore not a static monument; it is a bounded review surface whose claims
can strengthen or narrow as evidence changes.

\subsection{Minimal reader audit}

A reviewer who has only one hour should not start by reading every appendix.
The minimal audit is narrower. First, inspect the abstract and release delta to
see what kind of claim the release makes. Second, inspect the canonical tuple
and ask whether each component carries a distinct review question. Third, read
the K-level primer and ask whether levels are witness-bearing or merely named.
Fourth, inspect one theorem-bound claim and one replay-bound claim and check
whether the verbs match the support class. Fifth, inspect one comparator row and
ask whether OC has residual value after prior art is credited. If the release
survives those five checks, a deeper read is warranted. If it fails one of them,
the failure identifies the repair target without requiring the reviewer to
survey the whole manuscript.

\subsection{Why this discipline matters}

The manuscript is ambitious, and ambition invites a fair suspicion that the
language may be doing more work than the evidence. The evidence discipline in
this chapter is designed to answer that suspicion directly. It does not ask the
reader to trust the author's confidence. It asks the reader to inspect the
claim class, the support class, the comparator, and the reopening condition. In
that sense the publication surface is intentionally adversarial: a good
objection should find a named place to land, and a named place to land should
make the next repair possible.

\subsection{Boundary between manuscript and evidence files}

The manuscript explains the scientific argument. Evidence files preserve the
machine-readable details needed for replay, checksum comparison, and long-form
inspection. The two surfaces should not be confused. When a file name appears
in the manuscript, it is there as a stable public anchor, not as a substitute
for explanation. When a paragraph explains a theorem, replay, or comparator, it
must remain readable even before the reader opens the underlying file. This is
the standard used here: prose teaches the claim; evidence files make the claim
auditable.
